\chapter{Projeto 4: Meu Jogo Livre}

\begin{objetivos}
\begin{itemize}
\item Planejar um jogo criado por você do começo ao fim;
\item Testar, achar bugs e consertar;
\item Melhorar o projeto depois de testar e aprender como compartilhá-lo.
\end{itemize}
\end{objetivos}

\section{Agora é com você!}

Você já sabe tudo o que precisa: movimento, aparência, som, eventos, repetições, sensores, condições, variáveis, listas, clones e blocos próprios. Chegou a hora de criar o \textbf{seu} jogo!

\begin{missao}[Missão 20: Crie o seu jogo]
Invente, planeje, programe e teste um jogo ou uma história interativa. Se quiser, compartilhe o projeto depois.
\end{missao}

\section{Passo 1: escolha uma ideia}

Se estiver sem ideias, escolha uma destas:

\begin{itemize}
\item \textbf{Pega-pega:} um personagem foge de outro que anda atrás dele;
\item \textbf{Corrida:} dois jogadores apertam teclas diferentes e quem chegar primeiro vence;
\item \textbf{Nave espacial:} atire em inimigos que caem usando clones;
\item \textbf{Pulo do gato:} o gatinho pula obstáculos que andam na tela;
\item \textbf{História com escolhas:} o jogador responde perguntas e a história muda.
\end{itemize}

\section{Passo 2: planeje no papel}

Preencha este roteiro antes de arrastar o primeiro bloco:

\begin{definicao}[Meu roteiro]
\textbf{Nome do jogo:} \underline{\hspace{8cm}}\\[4pt]
\textbf{Atores:} \underline{\hspace{9cm}}\\[4pt]
\textbf{Cenários:} \underline{\hspace{9cm}}\\[4pt]
\textbf{Variáveis:} \underline{\hspace{9cm}}\\[4pt]
\textbf{Como se ganha:} \underline{\hspace{8cm}}\\[4pt]
\textbf{Como se perde:} \underline{\hspace{8cm}}
\end{definicao}

\section{Passo 3: programe aos poucos}

Não tente fazer tudo de uma vez! Siga esta ordem:

\begin{passos}
\item Faça o personagem principal \textbf{andar};
\item Adicione o \textbf{obstáculo} ou o \textbf{objetivo};
\item Crie as \textbf{regras} (pontos, vidas, condições);
\item Adicione \textbf{sons} e \textbf{efeitos};
\item Acrescente \textbf{tela inicial} e \textbf{fim de jogo}.
\end{passos}

\begin{dica}
Teste a cada pequeno passo! Clique na bandeira verde sempre que acrescentar um bloco novo. Assim, se algo der errado, você sabe exatamente onde procurar.
\end{dica}

\section{Passo 4: caçando os bugs}

\begin{vocabulario}[Vocabulário: bug e depuração]
\textbf{Bug} é um erro no programa que faz algo funcionar diferente do planejado. \textbf{Depurar} (ou \textbf{debugar}) é achar e consertar os bugs. Todo programador, mesmo os mais experientes, passa boa parte do tempo depurando!
\end{vocabulario}

Se o seu jogo estiver estranho, faça estas perguntas:

\begin{itemize}
\item Todos os blocos estão \textbf{encaixados}?
\item Cada pilha começa com um bloco de \textbf{evento}?
\item As variáveis são \textbf{zeradas} no começo do jogo?
\item O nome do ator ou da mensagem está \textbf{escrito igual} em todo lugar?
\item O ator está no \textbf{lugar} certo? Está escondido?
\end{itemize}

\begin{cuidado}[Bug dos bugs]
Quando nada funciona, \textbf{pare, respire} e execute o programa passo a passo, falando em voz alta o que cada bloco faz. Chama-se \textbf{depuração do patinho}: explicar o código para um patinho de borracha (ou para um amigo) ajuda a achar o erro!
\end{cuidado}

\section{Passo 5: compartilhe}

Quando o jogo estiver pronto, você pode compartilhá-lo no site do Scratch:

\begin{passos}
\item Crie uma conta gratuita em \texttt{scratch.mit.edu} (com a ajuda de um adulto);
\item Abra a mensagem enviada pelo Scratch e clique no link para confirmar o e-mail;
\item Dê um nome ao projeto no editor e clique em \textbf{Compartilhar};
\item Escreva em \textbf{Instruções} como se joga e em \textbf{Notas e créditos} quem fez o quê.
\end{passos}

\begin{curiosidade}
Na comunidade do Scratch você também pode \textbf{ver dentro} de milhões de projetos de outras pessoas e até \textbf{remixar}, isto é, copiar e melhorar com os seus próprios blocos, dando o crédito ao autor original!
\end{curiosidade}

\section{Hora de praticar}

\begin{exercicio}[Tela de início]
Adicione uma tela inicial ao seu jogo com o nome dele e o aviso ``Aperte espaço para começar''.
\end{exercicio}

\begin{solucao}
No Palco, use um cenário ``início'' e espere a tecla:
\begin{pilha}
\bk[scEvents]{quando \bandeira\ for clicado}
\bk[scLooks]{mude o cenário para \menu{início}}
\bk[scControl]{espere até que \booleano{tecla \menu{espaço} pressionada?}}
\bk[scLooks]{mude o cenário para \menu{jogo}}
\bk[scEvents]{envie \menu{começar}}
\end{pilha}
Os atores só começam o jogo quando recebem a mensagem \bloco[scEvents]{quando eu receber \menu{começar}}.
\end{solucao}

\begin{exercicio}[Recomeçar]
Faça o jogo recomeçar quando o jogador apertar a tecla R.
\end{exercicio}

\begin{solucao}
Use o evento de tecla para enviar a mesma mensagem da tela de início:
\begin{pilha}
\bk[scEvents]{quando a tecla \menu{r} for pressionada}
\bk[scEvents]{envie \menu{começar}}
\end{pilha}
No trecho da tela de início, mova a troca para o cenário \textbf{jogo} para uma pilha do Palco que recebe \textbf{começar}. Nessa pilha, zere também o placar. Crie a variável \textbf{pontos} se quiser testar este exemplo:
\begin{pilha}
\bk[scEvents]{quando eu receber \menu{começar}}
\bk[scLooks]{mude o cenário para \menu{jogo}}
\bk[scVariables]{defina \menu{pontos} para \entrada{0}}
\end{pilha}
Em cada ator, encaixe no receptor da mensagem os blocos que restauram a posição, a fantasia e a visibilidade inicial. Se o jogo usa vidas, tempo ou clones, restaure esses estados também. Ao encerrar uma rodada, deixe o receptor de R disponível. Se houver laços que começam com a mensagem, encerre os da rodada anterior antes de começar outra. Teste R durante e depois de uma partida e confira se a próxima começa do zero.
\end{solucao}

\begin{desafio}[O jogo final: teste e melhore]
\begin{passos}
\item Teste seu projeto no editor. Confira se a tela inicial aparece, se os controles funcionam, se o final esperado aparece e se dá para recomeçar. Se houver dois lados ou finais diferentes, teste um de cada vez.
\item Anote o que você esperava e o que aconteceu. Escolha uma coisa para melhorar, como a dica de um controle ou uma regra de reinício.
\item Faça a mudança e repita o mesmo teste. Compare o antes e o depois. Confira também se ainda dá para começar, terminar e jogar outra vez.
\end{passos}
\end{desafio}

\begin{solucao}
Exemplo da corrida. Ao apertar R por engano durante a disputa, os dois personagens voltavam à largada. A melhoria foi aceitar R só na tela final. Para conferir, mova os personagens e aperte R no meio da corrida: eles devem manter a posição. Termine a corrida e aperte R de novo: os dois devem voltar à largada. \textbf{Parabéns, programador!}
\end{solucao}
